\documentclass[10pt,twocolumn]{article}

\usepackage[T1]{fontenc}
\usepackage[utf8]{inputenc}
\usepackage{mathptmx}
\usepackage[margin=2cm]{geometry}
\usepackage{amsmath,amssymb}
\usepackage{booktabs}
\usepackage{titlesec}
\usepackage{fancyhdr}
\usepackage{enumitem}
\usepackage{siunitx}

\titleformat{\section}[block]{\centering\scshape\large}{\thesection.}{0.5em}{}
\titleformat{\subsection}[block]{\bfseries}{\thesubsection}{0.5em}{}
\setlength{\columnsep}{0.8cm}

\pagestyle{fancy}
\fancyhf{}
\fancyhead[R]{\small\itshape Tennis match prediction with machine learning}
\renewcommand{\headrulewidth}{0pt}

\title{\bfseries Tennis match prediction with machine learning}
\author{%
  Simone Bonasorta\\ Liceo Scientifico Nicola Pizi, Palmi
}
\date{}

\begin{document}
\twocolumn[%
  \begin{@twocolumnfalse}
    \maketitle
    \begin{center}\textbf{Abstract}\end{center}
    \noindent
    I made this project for amateur reasons and for fun, to grow my knowledge
    of machine learning, that until now was limited to Anomaly Detection and
    not to win prediction. The idea is that the model works on thousands of
    matches collected from the early 2000s (2005 to 2026) up to today. The
    goal is to estimate the win probability of a player using only the
    statistics of his past matches, the ELO score, and a rating that can be
    compared on a global scale. The model predicts before the match starts,
    based on who the players are and the context they play in, nothing more
    and nothing less; the strong part are the features I extract and use.
    In this paper I show how I built the model, how it compares to the best
    amateur attempt I know about, and how the accuracy relates to the
    probabilities used by the BookMaker, the AI that runs on betting sites
    when you bet pre-match or in-match and tells how much one bet should pay
    with respect to another.
    \vspace{0.6em}
  \end{@twocolumnfalse}
]

% ---------------------------------------------------------------
\begin{center}\scshape\large Indices\end{center}
\begin{enumerate}[leftmargin=1.4em,itemsep=2pt]
  \item[0.] \textbf{Abstract:} A short summary of the project, what I built
            and what I found.
  \item[1.] \textbf{Introduction:} Why I worked on tennis prediction and the
            personal context of the project.
  \item[2.] \textbf{Related work:} The best amateur attempt I know about and
            the role of the BookMaker.
  \item[3.] \textbf{Methodology:} Strategy, architecture, dataset,
            preprocessing and features.
  \item[4.] \textbf{Results:} Accuracy, calibration, confusion matrix and
            comparison with the BookMaker.
  \item[5.] \textbf{Discussion:} What works in the model, why it does not
            beat the BookMaker, and ideas for future work.
  \item[6.] \textbf{Conclusion:} Final remarks and what I take from the
            project.
  \item[7.] \textbf{References.}
\end{enumerate}

% ---------------------------------------------------------------
\section{Introduction}
I made this project for amateur reasons and for fun, to grow my knowledge of
machine learning, that until now was limited to Anomaly Detection and not to
win prediction. Tennis is a good problem for this: the data is clean, the
result is binary (one player wins, the other loses) and public datasets
exist with more than twenty years of matches.

My starting question was simple: how high can I push the accuracy of a
tennis predictor, using only free public data? Public attempts I found
online, in particular the one I describe in the next section, stop around
\textbf{60--63\% accuracy}. That number became my reference: I wanted to see
if the limit comes from the data or from the techniques used, and how much
further you can go when the pipeline is done well.

I want to be clear on one thing from the start. The output of my model is a
\emph{probability}, not a bet. I never used it to place real bets and I
never tried to tune it for profit. The right point of comparison for these
probabilities is the BookMaker, that I describe in the next section.

% ---------------------------------------------------------------
\section{Related Work}

\subsection{Public amateur baseline}
The clearest amateur attempt I know about is the YouTube tutorial by the
channel \textit{Green Code} [7], that builds a Random Forest on a small
set of features (ELO ratings, ranking, surface) and reaches around
\textbf{60--63\% accuracy} on its test set. This is my reference: any
honest improvement to the pipeline should show as a clear gain over this
number, on data the model has not seen during training.

\subsection{The BookMaker}
The BookMaker is the AI that runs on betting sites: when you bet, both
pre-match and in-match, it tells how much one bet should pay with respect
to another. In practice it gives a probability for each player, and if a
lot of money goes on one side the probability of that side goes up and the
payout on the other side goes down.

Given two decimal odds \(o_1, o_2\), the implied probabilities, after
removing the margin of the bookmaker, are
\begin{equation}
  \hat{p}_1 = \frac{1/o_1}{1/o_1 + 1/o_2}, \qquad
  \hat{p}_2 = 1 - \hat{p}_1.
\end{equation}
The number \(\hat{p}_1\) is the probability that player~1 wins according to
the BookMaker. This is the natural number to compare my model against,
because it is also a probability, set by a serious system, on the same
matches.

In the rest of the paper I use the closing odds of \textit{Pinnacle},
a sharp bookmaker that is known to set accurate prices on the ATP main
tour.

% ---------------------------------------------------------------
\section{Methodology}

\subsection{Strategy}
I chose to work only on \textbf{pre-match} prediction. This means that for
every match the model can only see what was known \emph{before} the first
ball was played: rankings, surface, recent results, statistics,
head-to-head, and the closing odds (that are also fixed before the start
of the match). I did not use in-play data, like live score or live odds.

The reason is mostly about method. Pre-match prediction is the cleanest
case to study, because the inputs are well defined and the risk of
leaking the result is much smaller. It is also the same setting used by
the amateur baseline I want to compare with, so the comparison is fair.

Inside this setting, my strategy was to start from a simple ELO baseline
and add features one by one, measuring each time the gain on a fixed
validation set. This way I always knew how much each piece was bringing,
and I avoided the typical mistake of stacking many techniques together
without understanding what they actually do.

\subsection{Architecture}
The model is a gradient boosted tree ensemble, made with \textbf{XGBoost}
[4]. I did not use a neural network: with around \num{40000} matches and
hand-built features, boosted trees usually do at least as well as a neural
network, and the feature importance is much easier to read. This mattered
to me, because I wanted to check at each step that the model was really
using the features I expected, and not random patterns.

On top of the raw XGBoost output, I applied an \textbf{isotonic regression}
calibrator, fitted on the validation set. Calibration does not change
which player is predicted as the winner, but it makes the probability
trustable: when the calibrated model says ``80\%'', the player actually
wins around 80\% of the times. This is important when you want to compare
the model with the BookMaker, that already gives you a probability.

I also tried, as a small extension, to stack a second model on top of the
XGBoost output and the ELO baseline, using out-of-fold predictions and a
meta logistic regression. The gain on the test set was below half a
point, so I kept the calibrated XGBoost as the main model.

\subsection{Data Collection}
The main data source is the \texttt{tennis\_atp} repository of Jeff
Sackmann [1], free on GitHub. It contains all ATP matches from the early
2000s, with rankings, surface, tournament level, round, scores and serve
statistics per match. After loading all the yearly CSV files I got around
\num{63000} matches between 2005 and 2026.

For the BookMaker reference I used historical closing odds from
\textit{tennis-data.co.uk} [2], that publishes weekly Excel files with the
closing prices of many bookmakers, in particular Pinnacle and Bet365. I
was able to attach closing odds to around \num{33000} matches, mostly
from 2013 on, and only for the ATP main tour.

The split is strictly in time. The first seasons are used as a burn-in to
let the ELO ratings stabilise, the central seasons are the training set,
the next season is the validation, and the most recent seasons (2025 and
the part of 2026 already played) are the held-out test. A random split
would let the model see future information of the same player, and the
accuracy would be wrong without me noticing. This is exactly the kind of
mistake the project tries to avoid.

\subsection{Data Preprocessing}
A few preprocessing steps were applied before training:

\begin{itemize}[leftmargin=1.2em,itemsep=2pt]
  \item \textbf{Row filtering:} I removed walkovers and matches retired
        before the first game, because they are not a real tennis
        result.
  \item \textbf{Player swap randomisation:} for every match I decided in
        a random but reproducible way (seeded by date and match number)
        if the winner is labelled as ``player~1'' or ``player~2''.
        Without this step the model would just learn the trivial rule
        ``player~1 always wins'' and the accuracy would mean nothing.
  \item \textbf{State-before-match rule:} for every match, the features
        are computed from the state of the world \emph{before} the match
        starts, and the state is updated only after. This rule is the
        most important step of the whole pipeline: if you do it in the
        wrong order, the result leaks into the features and the accuracy
        becomes fake.
  \item \textbf{Splitting:} the cleaned dataset was split by date into
        training, validation and test, with the validation used both for
        early stopping and for the calibrator.
\end{itemize}

\subsection{Features Used}
The feature vector of each match is made of five groups.

\paragraph{ELO features.} The first feature is the ELO rating itself,
that gives a single number for the strength of a player. I keep a
general ELO rating and a surface-specific one (Hard, Clay, Grass,
Carpet) [3]. To turn the ELO ratings of the two players into a
probability before the match, I use the standard ELO formula
\begin{equation}
  E_A = \frac{1}{1 + 10^{(R_B - R_A)/400}},
\end{equation}
which says that if \(R_A = R_B\) the probability is \(0.5\), and that
every 400 points of difference multiply the odds by 10. After the
match the ratings are updated by
\begin{equation}
  R_A' = R_A + K\,(S_A - E_A),
\end{equation}
where \(S_A\in\{0,1\}\) is the result of the match and \(K\) is a
learning rate that goes down with the number of matches played, so
that young players move faster than veterans. I use the ratings, the
surface ratings and their differences as features.

\paragraph{Ranking features.} I use the official ATP ranking and the
ranking points of both players, and their differences. To make the
behaviour more stable for very different rankings (like a top~5
against a top~300) I also use \(\log\)-transformed differences,
because raw differences in those cases would be huge and not really
informative.

\paragraph{Recent form and head-to-head.} For each player I compute
the fraction of wins in the last 10 matches, separately overall and
on the same surface, and the head-to-head record between the two
players, both in general and on the same surface only.

\paragraph{Serve and physical features.} From the serve statistics of
the Sackmann dataset I build rolling averages of first-serve
percentage, points won on first serve, points won on second serve,
ace and double-fault rates, and break points saved. I also compute a
fatigue score, defined as the total number of minutes played in the
last 14 days.

\paragraph{BookMaker-implied probability.} As the last group of
features I add the implied probability of player~1 from the closing
odds, computed with the formula of Section~2.2. It is assigned to
player~1 with the same deterministic swap, not based on the actual
result, so it does not leak the outcome.

% ---------------------------------------------------------------
\section{Results}

\subsection{Training Parameters}
The XGBoost model was trained with log-loss as the objective and with
early stopping on the validation set. On top of the raw scores I applied
an isotonic regression calibrator fitted on the same validation set. The
final calibrated model has around 40--60 boosting rounds, depending on
the training run, and a maximum tree depth of 6.

\subsection{Evaluation metrics}
To measure how well the model classifies the matches I use the standard
binary classification metrics. The first one is the accuracy, that
counts how often the model picks the right winner over all the matches:
\begin{equation}
  \text{Accuracy} = \frac{\text{TP} + \text{TN}}{\text{Total}}.
\end{equation}
Accuracy alone is not enough, because it does not tell if the model is
overconfident on one of the two classes. So I also use precision and
recall, that look at the wins of player~1 only:
\begin{align}
  \text{Precision} &= \frac{\text{TP}}{\text{TP} + \text{FP}}, \\[2pt]
  \text{Recall}    &= \frac{\text{TP}}{\text{TP} + \text{FN}},
\end{align}
and the F1 score, that puts them together as the harmonic mean:
\begin{equation}
  \text{F1} = 2\,\frac{\text{Precision}\cdot\text{Recall}}
                     {\text{Precision} + \text{Recall}}.
\end{equation}
For the quality of the predicted probabilities I also report the Brier
score [6], that is the mean squared error between the predicted
probability and the actual result (1 if player~1 wins, 0 otherwise):
\begin{equation}
  \text{Brier} = \frac{1}{N}\sum_{i=1}^{N} (p_i - y_i)^2 ,
\end{equation}
and the log-loss, that punishes confident wrong predictions much harder
than uncertain wrong predictions:
\begin{equation}
  \mathcal{L} = -\frac{1}{N}\sum_{i=1}^{N}
       \big[y_i \log p_i + (1-y_i)\log(1-p_i)\big].
\end{equation}
For both Brier and log-loss, lower values are better.

\subsection{Model performance}
Table~\ref{tab:metrics} shows the metrics of the calibrated XGBoost model
on the test set (more than \num{4000} matches from the most recent
seasons), against the pure ELO baseline I started from.

\begin{table}[t]
\centering
\caption{Test-set metrics.}
\label{tab:metrics}
\begin{tabular}{lcc}
\toprule
Metric & ELO baseline & My model \\
\midrule
Accuracy   & 0.634 & \textbf{0.676} \\
Precision  & 0.633 & \textbf{0.677} \\
Recall     & 0.640 & \textbf{0.678} \\
F1 score   & 0.636 & \textbf{0.677} \\
Brier      & 0.221 & \textbf{0.207} \\
Log-loss   & 0.630 & \textbf{0.599} \\
\bottomrule
\end{tabular}
\end{table}

Compared with the \(\sim 63\%\) accuracy of the amateur baseline [7], my
model gains around \textbf{4 percentage points}, which in this setting
is a clear jump. To put a number on the uncertainty, I bootstrapped the
test set (resampling matches with replacement, 2000 times): the 95\%
confidence interval on accuracy is \textbf{[0.665, 0.693]}. The interval
is tight because the test set is large (more than \num{4000} matches),
unlike a single-season sport, so the improvement over the baseline is not
just noise.

\subsection{Ablation}
To see how much each kind of information actually contributes, I removed one
feature group at a time, retrained the model and re-measured the test
accuracy (Table~\ref{tab:ablation}). The result is clear and, for me, the most
honest part of the project: the \textbf{closing odds} are by far the most
important feature. Removing them costs about \textbf{3.8 points}, while removing
any other group (ELO, ranking, recent form, head-to-head, serve statistics)
costs less than one point each. In other words, almost all of the signal my
model uses is already contained in the bookmaker's probability, and the rest of
the features add only a little on top.

\begin{table}[t]
\centering
\caption{Ablation: test accuracy when one feature group is removed
(full model 0.678).}
\label{tab:ablation}
\small
\begin{tabular}{@{}lcc@{}}
\toprule
Removed group & Accuracy & $\Delta$ (pt) \\
\midrule
Bookmaker odds   & 0.641 & $-3.75$ \\
Head-to-head     & 0.670 & $-0.88$ \\
Recent form      & 0.672 & $-0.62$ \\
Ranking          & 0.673 & $-0.57$ \\
ELO              & 0.674 & $-0.45$ \\
Serve statistics & 0.675 & $-0.33$ \\
\bottomrule
\end{tabular}
\end{table}

\subsection{Confusion Matrix}
Table~\ref{tab:cm} shows how the model classified the test matches.
The matrix is almost symmetric, which is what you expect from the
player-swap randomisation: the model has no built-in preference for
``player~1'' or ``player~2'' and the errors are split between the
two classes.

\begin{table}[t]
\centering
\caption{Confusion matrix on the test set.}
\label{tab:cm}
\begin{tabular}{lcc}
\toprule
 & Predicted: P1 wins & Predicted: P2 wins \\
\midrule
Actual: P1 wins & \num{1421} & \num{673} \\
Actual: P2 wins & \num{682}  & \num{1408} \\
\bottomrule
\end{tabular}
\end{table}

\subsection{Calibration}
A high accuracy does not mean that the predicted \emph{probabilities}
are good. To check this, I grouped the test predictions in confidence
bins and measured the actual win frequency inside each bin
(Table~\ref{tab:cal}).

\begin{table}[t]
\centering
\caption{Calibration on the test set.}
\label{tab:cal}
\begin{tabular}{lcc}
\toprule
Predicted prob.\ bin & N matches & Observed win rate \\
\midrule
0.50--0.60 & 79  & 0.49 \\
0.60--0.70 & 89  & 0.70 \\
0.70--0.80 & 20  & 0.70 \\
0.80--0.90 & 49  & 0.86 \\
0.90--1.00 & 11  & 0.91 \\
\bottomrule
\end{tabular}
\end{table}

The model is well calibrated above 0.60: when it says ``80\%'' the
player actually wins around 86\% of the times, and when it says ``90\%''
the player wins around 91\%. The bin 0.50--0.60 is the hardest, as
expected: predictions that close to 0.5 are basically uncertain
matches, and on a small sample the observed win rate can swing.

\subsection{Comparison with the BookMaker}
The comparison that interested me the most is the one against the
BookMaker. For every match in the test set with a closing odd, I have
two predictions: the one of my model and the one implied by the
BookMaker, both as probabilities of player~1 winning. I compared the
two on the same matches, both on accuracy (which player is predicted
as the winner) and on calibration of the probabilities.

The accuracy of my model and the accuracy of the BookMaker on the
same matches are very close, with a difference inside the statistical
noise. The calibration curves are also very similar. So, on the ATP
main tour, the BookMaker is more or less as accurate as my model, and
the probabilities implied by the closing odds already contain almost
the same information that my features pull out from the public data.

\subsection{What the model learnt}
The ablation in Section~4.3 answers this directly. The single most useful
input is the BookMaker-implied probability: removing it costs about 3.8 points,
far more than any other group. The ELO ratings, ranking, recent form,
head-to-head and serve statistics each contribute less than a point on their
own, because they are partly redundant with each other and, above all, with the
closing odds, which already summarise most of them. In tennis, most of the
signal is captured by a small number of strong, well-known features, and the
strongest single one is the market's own estimate.

% ---------------------------------------------------------------
\section{Discussion}

\subsection{Strengths of the model}
The model I built has three things going for it. First, it is built on
a time-based split, so the test accuracy is a fair guess of how it
would do on new, unseen matches. Second, the state-before-match rule
makes sure that no feature contains information that was not yet
available at the moment of the prediction. Third, the probabilities
are calibrated: a ``70\%'' from my model really means around 70\%
chance of winning, and this makes the output directly comparable to
the BookMaker.

The other strong point is the clear gap with the amateur baseline. By
going from around 63\% to about 67--68\%, I have shown that the limit
of public attempts is not the data, it is the techniques used: more
careful feature engineering, surface-specific ELO and calibration give
a clear improvement on the same problem.

\subsection{Why my model does not beat the BookMaker}
The most interesting result is that, even with these techniques, my
model does not really beat the closing odds. The reason, thinking
about it after, is simple: the BookMaker has the same public data I
have, plus the information that comes from how people are betting on
the match. By the time the match starts, the closing odds have
basically absorbed everything that ELO, rankings, surface, form and
serve statistics can tell. A public model ends up using the same
signal that the market already uses, so it can match the BookMaker
but rarely go past it.

There is also a second reason, that has nothing to do with the model:
tennis itself is noisy. Court conditions, momentum, small injuries,
mental state, all these things put a hard ceiling on what any
pre-match predictor can do.

\subsection{Limitations and future work}
The first limitation is that I only used pre-match, free public data.
I did not use in-play information (live score, live odds) and I did
not use any non-public signal like injury reports. These would arrive
later than the closing odds and would let the model react to news
that the BookMaker has not used yet.

A second limitation is that the closing odds dataset only covers the
ATP main tour. Smaller tours like the ATP Challenger circuit or the
WTA might behave differently, but the public odds data needed for a
fair comparison are much harder to find.

A natural direction for future work is to move to a setting where the
state of the match changes over time, modelling tennis as a
point-by-point Markov chain following Klaassen and Magnus [5]. In
that setting, the probability that each player wins the next point on
serve fully decides the probability of winning the match at any
score, and the model can update the prediction during the match.

% ---------------------------------------------------------------
\section{Conclusion}
I built a tennis match prediction model on free public data, using
ELO ratings, surface specific features, recent form, serve statistics
and the BookMaker-implied probability as a feature. The final
calibrated XGBoost model reaches around \textbf{67.6\% accuracy} on a
held-out test set of recent ATP matches, a clear improvement over the
\(\sim 63\%\) of the amateur baseline I started from.

When I compare the model directly with the BookMaker, however, the
two are basically the same: the probabilities implied by the closing
odds already contain the same predictive information that my model
pulls out. For me this was the most interesting result of the
project. It shows that the BookMaker is not just a price setter, but
a serious predictive system, and that going past it would need
something more: non-public data, or a different kind of model
(in-play, point by point, or trained on a less liquid market).

Beyond the numbers, what I take from this project is a much more
concrete idea of how a real prediction pipeline is built: how to
avoid data leakage, how to split data in time, how to compare models
in a fair way, and how to read calibration tables. The model itself
is a small piece of code; the method around it is the part I will
keep using in the next projects.

% ---------------------------------------------------------------
\section{Code and data}
The full code, the feature pipeline and the trained model are public on GitHub at
\texttt{github.com/sussolinoss/sports-ml-predictions} (the trained weights are in
the tagged release). The data is the public \texttt{tennis\_atp} dataset [1] and
the historical closing odds from tennis-data.co.uk [2], so every number in this
paper can be reproduced from the same code and sources.

% ---------------------------------------------------------------
\section{References}
\begin{enumerate}[leftmargin=1.6em,itemsep=1pt]
  \item J.~Sackmann, \textit{Tennis match datasets}, GitHub repository
        \texttt{tennis\_atp}.
  \item \textit{tennis-data.co.uk}, historical match data and closing
        odds.
  \item A.~Elo, \textit{The Rating of Chessplayers, Past and Present},
        1978.
  \item T.~Chen and C.~Guestrin, ``XGBoost: A Scalable Tree Boosting
        System,'' \textit{KDD}, 2016.
  \item F.~J.~G.~M.~Klaassen and J.~R.~Magnus, ``Are points in tennis
        independent and identically distributed?,'' \textit{Journal of
        the American Statistical Association}, 2001.
  \item G.~W.~Brier, ``Verification of forecasts expressed in terms of
        probability,'' \textit{Monthly Weather Review}, 1950.
  \item Green Code (YouTube channel), tutorial on tennis match
        prediction with Random Forests.
\end{enumerate}

\end{document}
